% !TEX program = xelatex
\documentclass[a4paper,14pt,oneside,openany]{memoir}

%%% Задаем поля, отступы и межстрочный интервал %%%

\usepackage[left=30mm, right=15mm, top=20mm, bottom=20mm]{geometry}
\pagestyle{plain}
\parindent=1.25cm
\usepackage{indentfirst}

%%% Задаем языковые параметры и шрифт %%%

\usepackage[english,russian]{babel}
\babelfont[russian]{rm}{Times New Roman}
\babelfont[russian]{tt}{Courier New}

%%% Задаем стиль заголовков %%%

\setsecnumdepth{subsection}
\renewcommand*{\chapterheadstart}{}
\renewcommand*{\printchaptername}{}
\renewcommand*{\chapnumfont}{\normalfont\bfseries}
\renewcommand*{\afterchapternum}{\hspace{1em}}
\renewcommand*{\printchaptertitle}{\normalfont\bfseries\centering\MakeUppercase}
\setbeforesecskip{20pt}
\setaftersecskip{20pt}
\setsecheadstyle{\raggedright\normalfont\bfseries}
\setbeforesubsecskip{20pt}
\setaftersubsecskip{20pt}
\setsubsecheadstyle{\raggedright\normalfont\bfseries}

\setsecnumdepth{none}
\counterwithout{figure}{chapter}
\counterwithout{table}{chapter}

%%% Параметры переносов %%%

\tolerance 1414
\hbadness 1414
\emergencystretch 1.5em
\clubpenalty=10000
\widowpenalty=10000

%%% Параметры рисунков и листингов %%%

\usepackage{graphicx}
\usepackage{caption}
\usepackage{subcaption}
\usepackage{float}
\usepackage{listings}
\usepackage{xcolor}

\captionsetup[figure]{font=small, width=\textwidth, name=Рисунок, justification=centering}
\captionsetup[table]{font=small, width=\textwidth, name=Таблица, justification=raggedright}
\captiondelim{ --- }

%%% Настройка отображения листингов кода %%%
\definecolor{codebg}{rgb}{0.97,0.97,0.97}
\definecolor{commentgreen}{rgb}{0,0.4,0}
\lstset{
    backgroundcolor=\color{codebg},
    basicstyle=\ttfamily\small,
    breaklines=true,
    commentstyle=\color{commentgreen},
    keywordstyle=\color{blue},
    numbers=left,
    numberstyle=\tiny\color{gray},
    numbersep=8pt,
    frame=single,
    frameround=tttt,
    showstringspaces=false,
    extendedchars=true,
    inputencoding=utf8
}
\renewcommand{\lstlistingname}{Листинг}

%%% Гиперссылки %%%

\usepackage{hyperref}
\hypersetup{
    colorlinks=true,
    linkcolor=blue,
    citecolor=blue,
    urlcolor=blue,
    breaklinks=true
}

%%% Списки %%%

\usepackage{enumitem}
\renewcommand*{\labelitemi}{\normalfont{--}}
\setlist{noitemsep, leftmargin=*}

\begin{document}

% ==========================================
% ТИТУЛЬНЫЙ ЛИСТ
% ==========================================
\thispagestyle{empty}

\begin{center}
    МИНИСТЕРСТВО ЦИФРОВОГО РАЗВИТИЯ, СВЯЗИ И МАССОВЫХ \\ КОММУНИКАЦИЙ \\ РОССИЙСКОЙ ФЕДЕРАЦИИ

    \vspace{20pt}

    Федеральное государственное бюджетное образовательное учреждение \\ высшего образования \\
    «Сибирский государственный университет телекоммуникаций и \\ информатики»

    \vspace{20pt}

    Кафедра телекоммуникационных систем и вычислительных средств \\ (ТСиВС)
\end{center}

\vfill

\begin{center}
    Расчетно-графическое задание \\
    по дисциплине \\
    \textit{«WEB-Технологии»}

    \vspace{20pt}

    по теме: \\
    \uppercase{Создание тестовой видеоплатформы с использованием технологий Python для бэкенда, React JS для фронтенда и Figma для дизайна}
\end{center}

\vfill

\noindent Студент: \\
\textit{Группа № ИКС-532 \hfill Лёшин А. С.}

\vspace{20pt}

\noindent Преподаватель: \\
\textit{Старший преподаватель \hfill Андреев А.В.}

\vfill

\begin{center}
    Новосибирск 2026 г.
\end{center}

\newpage
\setcounter{page}{2}
\OnehalfSpacing*

% ==========================================
% 1. ВВЕДЕНИЕ
% ==========================================
\section{1 Введение}

Проект представляет собой современное веб-приложение с распределенной архитектурой клиент-сервер, использующее фреймворк Django для реализации серверного бэкенда и библиотеку React для построения интерактивного фронтенда.

\subsection{1.1 Задачи проекта}
\begin{itemize}
    \item Разработка структурированного и документированного REST API на базе Django REST Framework.
    \item Создание надежной системы регистрации и авторизации пользователей с использованием безопасных JWT-токенов.
    \item Реализация серверного и клиентского функционала для загрузки, обработки и потокового воспроизведения медиафайлов.
    \item Разработка отзывчивого пользовательского интерфейса на React с применением утилитарного CSS-фреймворка Tailwind CSS.
    \item Конфигурация серверного окружения и деплой готового проекта на удаленный VDS сервер.
\end{itemize}

% ==========================================
% 2. АРХИТЕКТУРА ПРОЕКТА
% ==========================================
\section{2 Архитектура проекта}

\subsection{2.1 Стек технологий}
Компоненты разработанной системы и их технологическое назначение подробно систематизированы и представлены в таблице 1.

\begin{table}[H]
\centering
\small
\caption{Технологическая спецификация компонентов}
\begin{tabular}{|l|l|}
\hline
\textbf{Компонент системы} & \textbf{Используемая технология} \\ \hline
Бэкенд-платформа & Django 5.0, Django REST Framework \\ \hline
Аутентификация сессий & JWT (djangorestframework-simplejwt) \\ \hline
База данных & Реляционная СУБД SQLite 3 \\ \hline
Фронтенд-архитектура & React 18, React Router DOM v6 \\ \hline
Стилизация интерфейса & Tailwind CSS \\ \hline
Асинхронный транспорт & HTTP-клиент Axios \\ \hline
Веб-сервер удаленного деплоя & Nginx (Amvera Cloud / VDS) \\ \hline
\end{tabular}
\end{table}

\subsection{2.2 Структура проекта}
Разделение кодовой базы на независимые каталоги бэкенда и фронтенда обеспечивает изолированную разработку каждого слоя системы. Логическое дерево каталогов проекта приведено в листинге 1.

\begin{lstlisting}[caption={Структура проекта MyTube}]
video-platform/
 backend/             # Django бэкенд
    core/            # Настройки проекта
    users/           # Приложение управления пользователями
    videos/          # Приложение управления видеоконтентом
    media/           # Хранилище загруженных файлов
 frontend/            # React фронтенд
     public/          # Статические файлы
     src/
         components/  # Переиспользуемые React компоненты
         pages/       # Страницы приложения
         contexts/    # Глобальные стейты (Context API)
         App.js       # Главный корневой компонент
\end{lstlisting}

% ==========================================
% 3. РАЗРАБОТКА БЭКЕНДА
% ==========================================
\section{3 Разработка бэкенда}

\subsection{3.1 Модели данных}
Архитектура базы данных построена на расширении стандартной модели пользователя для кастомизации профилей (листинг 2) и сущности хранения метаданных видеоконтента (листинг 3).

\begin{lstlisting}[language=Python, caption={Модель пользователя User}]
class User(AbstractUser):
    avatar = models.ImageField(upload_to='avatars/', null=True, blank=True)
    banner = models.ImageField(upload_to='banners/', null=True, blank=True)
    bio = models.TextField(max_length=500, blank=True, default='')
    location = models.CharField(max_length=100, blank=True, default='')
    website = models.URLField(blank=True, default='')
    birth_date = models.DateField(null=True, blank=True)
    subscribers = models.ManyToManyField('self', symmetrical=False, blank=True)
\end{lstlisting}

\begin{lstlisting}[language=Python, caption={Модель видео Video}]
class Video(models.Model):
    title = models.CharField(max_length=200)
    description = models.TextField(blank=True)
    video_file = models.FileField(upload_to='videos/%Y/%m/')
    thumbnail = models.ImageField(upload_to='thumbnails/%Y/%m/', null=True, blank=True)
    author = models.ForeignKey(User, on_delete=models.CASCADE, related_name='videos')
    views = models.IntegerField(default=0)
    likes = models.ManyToManyField(User, related_name='liked_videos', blank=True)
    created_at = models.DateTimeField(auto_now_add=True)
\end{lstlisting}

\subsection{3.2 API Эндпоинты}
Маршрутизация запросов к REST API серверной части представлена в таблице 2.

\begin{table}[H]
\centering
\small
\caption{Маршруты REST API подсистемы}
\begin{tabular}{|l|l|l|}
\hline
\textbf{Метод} & \textbf{Эндпоинт} & \textbf{Описание функционала} \\ \hline
\textbf{POST} & /api/users/register/ & Регистрация нового аккаунта \\ \hline
\textbf{POST} & /api/users/login/ & Вход в систему (генерация JWT токена) \\ \hline
\textbf{GET} & /api/users/me/ & Получение данных текущей сессии \\ \hline
\textbf{GET} & /api/videos/ & Получение полного списка видеороликов \\ \hline
\textbf{GET} & /api/videos/recommended/ & Запрос блока рекомендованных видео \\ \hline
\textbf{POST} & /api/videos/ & Загрузка нового видео на сервер \\ \hline
\textbf{POST} & /api/videos/\{id\}/like/ & Переключение лайка на контенте \\ \hline
\textbf{POST} & /api/videos/\{id\}/view/ & Инкремент счетчика просмотров \\ \hline
\textbf{GET} & /api/videos/\{id\}/comments/ & Запрос комментариев под видео \\ \hline
\textbf{POST} & /api/videos/\{id\}/comments/ & Добавление текстового комментария \\ \hline
\textbf{DELETE} & /api/videos/\{id\}/ & Изолированное удаление видеоролика \\ \hline
\textbf{PATCH} & /api/users/\{id\}/update-profile/ & Валидация и обновление профиля \\ \hline
\end{tabular}
\end{table}

\subsection{3.3 JWT Аутентификация}
Интеграция авторизационных ограничений в настройки проекта Django представлена в листинге 4.

\begin{lstlisting}[language=Python, caption={Настройка JWT в файле settings.py}]
REST_FRAMEWORK = {
    'DEFAULT_AUTHENTICATION_CLASSES': (
        'rest_framework_simplejwt.authentication.JWTAuthentication',
    ),
}

SIMPLE_JWT = {
    'ACCESS_TOKEN_LIFETIME': timedelta(days=1),
    'REFRESH_TOKEN_LIFETIME': timedelta(days=7),
}
\end{lstlisting}

% ==========================================
% 4. РАЗРАБОТКА ФРОНТЕНДА
% ==========================================
\section{4 Разработка фронтенда}

\subsection{4.1 Основные компоненты}

\subsubsection{4.1.1 VideoCard --- карточка отображения видео}
Компонент отвечает за рендеринг превью, заголовка и автора контента в общей сетке выдачи (листинг 5).

\begin{lstlisting}[language=Java, caption={Компонент карточки видео VideoCard.jsx}]
export default function VideoCard({ video }) {
  return (
    <Link to={`/video/${video.id}`} className="block group">
      <div className="bg-gray-900 rounded-xl overflow-hidden">
        <div className="relative pb-[56.25%]">
          {video.thumbnail ? (
            <img src={video.thumbnail} alt={video.title} className="absolute inset-0 w-full h-full object-cover" />
          ) : (
            <div className="absolute inset-0 bg-gradient-to-br from-purple-900 to-pink-900" />
          )}
        </div>
        <div className="p-3">
          <h3 className="font-semibold text-white">{video.title}</h3>
          <p className="text-sm text-gray-400">{video.author?.username}</p>
        </div>
      </div>
    </Link>
  );
}
\end{lstlisting}

\subsubsection{4.1.2 Механизм динамического поиска}
Запросы к API-серверу для фильтрации контента по строковым паттернам без перезагрузки страницы представлены в листинге 6.

\begin{lstlisting}[language=Java, caption={Реализация поисковой фильтрации}]
const performSearch = async (query) => {
  const response = await axios.get(`/videos/?search=${encodeURIComponent(query)}`);
  setVideos(response.data.results || response.data);
};

useEffect(() => {
  const params = new URLSearchParams(window.location.search);
  const query = params.get('search');
  if (query) performSearch(query);
}, [window.location.search]);
\end{lstlisting}

\subsection{4.2 AuthContext --- управление глобальным стейтом авторизации}
Использование React Context API для обеспечения сквозной аутентификации сессий приведено в листинге 7.

\begin{lstlisting}[language=Java, caption={Контекст авторизации AuthContext.jsx}]
const AuthContext = createContext();

export const AuthProvider = ({ children }) => {
  const [user, setUser] = useState(null);

  const login = async (username, password) => {
    const response = await axios.post('/users/login/', { username, password });
    localStorage.setItem('access_token', response.data.access);
    setUser(response.data.user);
    return true;
  };

  return (
    <AuthContext.Provider value={{ user, login, logout }}>
      {children}
    </AuthContext.Provider>
  );
};
\end{lstlisting}

% ==========================================
% 5. ДЕПЛОЙ НА СЕРВЕР
% ==========================================
\section{5 Деплой на сервер}

\subsection{5.1 Настройка системного окружения}
Установка необходимого серверного программного обеспечения на Linux-дистрибутив VDS (листинг 8).

\begin{lstlisting}[language=bash, caption={Установка зависимостей на сервере}]
apt update && apt upgrade -y
apt install -y python3-pip python3-venv nginx nodejs npm
\end{lstlisting}

\subsection{5.2 Запуск бэкенда через Gunicorn}
Создание изолированного окружения и инициализация многопоточного WSGI-сервера (листинг 9).

\begin{lstlisting}[language=bash, caption={Настройка и старт процессов Gunicorn}]
cd /root/video-platform/backend
python3 -m venv venv
source venv/bin/activate
pip install -r requirements.txt
gunicorn --workers 3 --bind 0.0.0.0:8000 backend.wsgi:application
\end{lstlisting}

\subsection{5.3 Конфигурация веб-сервера Nginx}
Настройка проксирования входящих запросов к API и раздача статического билда фронтенда (листинг 10).

\begin{lstlisting}[language=bash, caption={Конфигурационный файл Nginx}]
server {
    listen 80;
    server_name 62.109.1.29;

    location /api/ {
        proxy_pass http://127.0.0.1:8000/api/;
    }

    location / {
        root /root/video-platform/frontend/build;
        try_files $uri $uri/ /index.html;
    }
}
\end{lstlisting}

% ==========================================
% 7. РЕЗУЛЬТАТЫ РАБОТЫ
% ==========================================
\section{7 Результаты работы}
В ходе выполнения проекта были достигнуты следующие практические результаты:
\begin{itemize}
    \item Разработана и отлажена безопасная регистрация и авторизация пользователей с JWT-токенами.
    \item Реализован плеер для плавного просмотра видеоконтента.
    \item Реализован асинхронный модуль загрузки видеоматериалов с генерацией уникальных превью.
    \item Интегрирована интерактивная система лайков и текстовых комментариев.
    \item Разработан модуль подписки на каналы выбранных авторов.
    \item Внедрен реактивный поиск по названиям видеороликов на главном экране.
    \item Обеспечено редактирование профилей (кастомизация баннера, аватара, биографии).
\end{itemize}

% ==========================================
% 8. ВЫВОДЫ
% ==========================================
\section{8 Выводы}
В ходе выполнения расчетно-графического задания была успешно спроектирована и программно реализована полноценная видеоплатформа «MyTube». Разработано отказоустойчивое REST API средствами Django REST Framework с JWT-аутентификацией сессий. Создан интерактивный визуальный фронтенд-компонент на React с применением утилитарной стилизации Tailwind CSS. Все ключевые модули, включая подписки, комментарии, динамические счетчики и загрузку мультимедиа, были успешно интегрированы и развернуты на удаленном сервере VDS. Поставленные задачи выполнены в полном объеме, закреплен практический опыт построения современных клиент-серверных приложений.

% ==========================================
% 9. ССЫЛКИ НА РЕСУРСЫ
% ==========================================
\section{9 Ссылки на ресурсы}
\begin{itemize}
    \item \textbf{GitHub репозиторий проекта:} \\ \url{https://github.com/ArtemLeshin/video-hosting}
    \item \textbf{Адрес опубликованного веб-интерфейса (Netlify):} \\ \url{https://exquisite-dusk-620e3d.netlify.app}
    \item \textbf{Базовый адрес бэкенд-сервера API (Amvera):} \\ \url{https://mutube-dreamshelter.amvera.io}
    \item \textbf{Презентация выполнения работы сайта находится в самом низу файла README.md} \\ \url{https://github.com/ArtemLeshin/video-hosting}
\end{itemize}

\end{document}